\begin{lemma}[Quantum group characters at $q=e^{i\pi/(k+2)}$]
\label{P14:lem:qgroup_char}
For $\mathrm{SU}(2)_k$ at $q=e^{i\pi/(k+2)}$:
\begin{equation}
  d_j = \chi_j(q) = \sum_{m=-j}^{j}q^{2m}
  = \frac{\sin((2j+1)\pi/(k+2))}{\sin(\pi/(k+2))}.
  \label{P14:eq:qchar}
\end{equation}
At $k=3$: $d_0=1$, $d_{1/2}=\vp$, $d_1=\vp$, $d_{3/2}=1$,
and $\vp = 1+2\cos(2\pi/5)$ is the character of the $j=1$ representation.
The truncation $[5]_q=0$ implements the pentagon truncation.
\status{published}
\residual{exact; $\vp=1+2\cos(2\pi/5)$ is the $j=1$ quantum group character at $q=e^{i\pi/5}$}
\source{P14:lem:qgroup_char}
\end{lemma}
\status{published}
\source{P14:lem:qgroup_char}

\begin{lemma}[Parseval identity for the WZW primary spectrum]
\label{P14:lem:parseval}
\begin{equation}
  \sum_{j=0}^{k/2}\sin^2\!\frac{(2j+1)\pi}{k+2} = \frac{k+2}{2}.
  \label{P14:eq:parseval_wzw}
\end{equation}
For $k+2=5$ (pentagon): sum $= 5/2 = (\text{pentagon sides})/2$.
Proof: $\sum_{n=1}^{k+1}\cos(2n\pi/(k+2)) = \operatorname{Re}[-1]=-1$
from the complete $(k+2)$-th roots of unity summing to zero.
\status{published}
\residual{exact; holds for any $N=k+2$; the result $5/2$ is the pentagon's Parseval sum}
\source{P14:lem:parseval}
\end{lemma}
\status{published}
\source{P14:lem:parseval}

\begin{corollary}[Total quantum dimension identity]
\label{P14:cor:D2_parseval}
\begin{equation}
  D^2\,\sin^2\!\frac{\pi}{k+2} = \frac{k+2}{2},
  \label{P14:eq:D2_sin2}
\end{equation}
where $D^2=\sum_j d_j^2 = 5+\sqrt{5}$ at $k=3$.
Check: $(5+\sqrt{5})(5-\sqrt{5})/8 = 20/8 = 5/2$. $\checkmark$
\status{published}
\source{P14:cor:D2_parseval}
\end{corollary}
\status{published}
\source{P14:cor:D2_parseval}

\begin{lemma}[KZB monodromy and torus sector weights]
\label{P14:lem:kzb_monodromy}
For $\mathrm{SU}(2)_k$ on the torus $T^2$ with A-type modular invariant,
the $b$-cycle monodromy of the torus conformal blocks is the modular
S-matrix $S_{jl}$ (Falceto--Gaw\c{e}dzki~1996; Felder--Wieczerkowski~1996).
The four-point contact term in
$Z^{-1}\mathrm{Tr}[q^{L_0-c/24}\phi_{1/2}^4(0)]$
receives per-sector contribution
\begin{equation}
  w_j^{\mathrm{torus}} = |S_{0j}|^2\cdot\frac{k+2}{2}
  = d_j^2\,\sin^2\!\frac{\pi}{k+2},
  \label{P14:eq:kzb_weight}
\end{equation}
replacing the plane OPE weight
$w_j^{\mathrm{plane}} = (C^j_{1/2,1/2})^2\,d_j$.
\status{proved (cites Falceto--Gaw\c{e}dzki 1996)}
\residual{this is the step that closes the proof of Theorem~\ref{P14:thm:higgs}; without it the result is a Proposition}
\source{P14:lem:kzb_monodromy; Falceto--Gaw\c{e}dzki (1996)}
\end{lemma}
\status{published}
\source{P14:lem:kzb_monodromy}

\begin{theorem}[Pentagon identity and the torus four-point contact term]
\label{P14:thm:pentagon_g4}
\begin{equation}
  G_4^{\mathrm{torus}}[\phi_{1/2}^4]
  = \sum_{j}d_j^2\,\sin^2\!\frac{\pi}{k+2}
  = D^2\,\sin^2\!\frac{\pi}{k+2}
  = \frac{k+2}{2}.
  \label{P14:eq:G4_pentagon}
\end{equation}
Torus-to-plane enhancement: $G_4^{\mathrm{torus}}/G_4^{\mathrm{plane}}
= (k+2)/2\,/\,(3-\vp) = \sqrt{5}\vp/2 = (5+\sqrt{5})/4 \approx 1.809$.
\status{proved (using Lemmas~\ref{P14:lem:kzb_monodromy}, \ref{P14:lem:parseval}, \ref{P14:lem:qgroup_char})}
\residual{$G_4^{\mathrm{torus}} = 5/2$ exactly; $(k+2)/2 = (\text{pentagon sides})/2$}
\source{P14:thm:pentagon_g4}
\end{theorem}
\status{published}
\source{P14:thm:pentagon_g4}

\begin{remark}[The pentagon gives $G_4^{\mathrm{torus}}$ directly]
\label{P14:rem:pentagon_direct}
The identity $G_4^{\mathrm{torus}}=(k+2)/2$ results from two standard identities
tracing to the same root: $k+2=5$ is both the number of pentagon sides
and the order of the root of unity $q=e^{i\pi/5}$ at which the quantum group truncates.
The two-step identity
\begin{equation}
  G_4^{\mathrm{torus}}
  = D^2\,\sin^2\!\frac{\pi}{k+2}
  = \sum_{j=0}^{k/2}\sin^2\!\frac{(2j+1)\pi}{k+2}
  = \frac{k+2}{2}
\end{equation}
requires: (i) characters $d_j=\chi_j(e^{i\pi/5})$ (Lemma~\ref{P14:lem:qgroup_char});
(ii) the Parseval sum for $N=k+2=5$ (Lemma~\ref{P14:lem:parseval});
(iii) the per-sector modular weight $\sin^2(\pi/(k+2))$.
The result $5/2=(\text{pentagon sides})/2$ is the Parseval sum for $N=5$.
\status{published}
\source{P14:rem:pentagon_direct}
\end{remark}
\status{published}
\source{P14:rem:pentagon_direct}

\begin{remark}[Modular renormalisation of OPE weights]
\label{P14:rem:modular_renorm}
The plane-to-torus change is a change of per-sector OPE weight:
\begin{align}
  \text{plane:}&\quad w_j = (C^j_{1/2,\,1/2})^2\,d_j\quad\text{(OPE coefficients squared)},\notag\\
  \text{torus:}&\quad w_j = d_j^2\,\sin^2\!\frac{\pi}{k+2} = |S_{0j}|^2\cdot\frac{k+2}{2}\quad\text{(modular weights)}.
\end{align}
Ratio: $G_4^{\mathrm{torus}}/G_4^{\mathrm{plane}} = D^2\sin^2(\pi/(k+2))/(3-\vp)
= (k+2)/2\,/\,(5-\sqrt{5})/2 = \sqrt{5}\vp/2$.
The weight change is established by the $b$-cycle KZB monodromy (Lemma~\ref{P14:lem:kzb_monodromy}),
which together with the Parseval identity (Lemma~\ref{P14:lem:parseval}) completes the proof
that $G_4^{\mathrm{torus}}=(k+2)/2$.
\status{published}
\source{P14:rem:modular_renorm}
\end{remark}
\status{published}
\source{P14:rem:modular_renorm}

\begin{lemma}[Plane quartic coupling from the Fibonacci $F$-matrix]
\label{P14:lem:fmatrix_plane}
In $\mathrm{SU}(2)_3$ on the plane, $(C^0)^2 = 1/\vp^2$,
$(C^1)^2 = 1/\vp$ (Fibonacci $F$-matrix), giving:
\begin{equation}
  G_4^{\mathrm{plane}} = (C^0)^2 d_0 + (C^1)^2 d_1
  = 1/\vp^2 + 1 = 3-\vp.
  \label{P14:eq:G4plane}
\end{equation}
\begin{equation}
  \lambda_{\mathrm{plane}} = (3-\vp)\,r_{\mathrm{WZW}}^2
  \quad\Rightarrow\quad
  m_H^{\mathrm{plane}} \approx 93\,\mathrm{GeV}.
  \label{P14:eq:lambda_plane}
\end{equation}
\status{published}
\residual{plane result inconsistent with observation; torus enhancement (Theorem~\ref{P14:thm:pentagon_g4}) required}
\source{P14:lem:fmatrix_plane}
\end{lemma}
\status{published}
\source{P14:lem:fmatrix_plane}

\begin{theorem}[Algebraic closure of the Higgs coupling]
\label{P14:thm:higgs_algebra}
\begin{equation}
  \frac{k+2}{2} = (3-\vp)\cdot\frac{\sqrt{5}\,\vp}{2}.
  \label{P14:eq:higgs_closure}
\end{equation}
Proof: $(3-\vp)\cdot\sqrt{5}\vp/2 = \sqrt{5}(2\vp-1)/2 = \sqrt{5}\cdot\sqrt{5}/2 = 5/2$.
Uses only $\vp^2=\vp+1$ and $2\vp-1=\sqrt{5}$.
\status{proved (algebra)}
\source{P14:thm:higgs_algebra}
\end{theorem}
\status{published}
\source{P14:thm:higgs_algebra}

\begin{remark}[Torus enhancement and proof completion]
\label{P14:rem:torus_enhancement}
The enhancement factor $\sqrt{5}\vp/2=(5+\sqrt{5})/4\approx1.809$
is established by Lemma~\ref{P14:lem:kzb_monodromy} (Falceto--Gaw\c{e}dzki~1996):
the $b$-cycle KZB monodromy is the modular S-matrix, replacing plane OPE weights
$(C^j)^2 d_j$ with torus modular weights $d_j^2\sin^2(\pi/(k+2))$,
summing to $(k+2)/2=5/2$ (Corollary~\ref{P14:cor:D2_parseval}).
The algebraic form connects to the Verlinde one-point function:
$\langle\phi_{1/2}\rangle_{\mathrm{torus}} = S_{0,1/2}/S_{0,0}=\vp$.
The enhancement $\sqrt{5}\vp/2=\vp\cdot\sqrt{5}/2$ is the product of this
Verlinde VEV $\vp$ and the pentagon factor $\sqrt{5}/2$.
The plane result $3-\vp$ recovers the torus result via Theorem~\ref{P14:thm:higgs_algebra}:
$(3-\vp)\cdot\vp\sqrt{5}/2=5/2$.
\status{published}
\source{P14:rem:torus_enhancement}
\end{remark}
\status{published}
\source{P14:rem:torus_enhancement}

\begin{theorem}[Higgs mass from the WZW condensate]
\label{P14:thm:higgs}
In the condensate framework at WZW level $k=3$:
\begin{equation}
  \boxed{\lambda = \frac{k+2}{2}\,r_{\mathrm{WZW}}^2
  = \frac{5\cdot 2^{5/3}}{\vp^{10}} = 0.12907}
  \label{P14:eq:lambda_higgs}
\end{equation}
\begin{equation}
  \boxed{m_H = \sqrt{k+2}\;r_{\mathrm{WZW}}\;\vEW
  = \sqrt{5}\;\frac{2^{4/3}}{\vp^5}\;\vEW = 125.10\,\mathrm{GeV}}
  \label{P14:eq:mH_formula}
\end{equation}
Proof path (five items, all proved):
\begin{enumerate}[noitemsep,label=(\roman*)]
\item $d_j=\chi_j(e^{i\pi/5})$; $\vp=1+2\cos(2\pi/5)$ (Lemma~\ref{P14:lem:qgroup_char})
\item Parseval: $\sum_j\sin^2((2j+1)\pi/(k+2))=(k+2)/2$ (Lemma~\ref{P14:lem:parseval})
\item $D^2\sin^2(\pi/(k+2))=(k+2)/2$ (Corollary~\ref{P14:cor:D2_parseval})
\item $\lambda_{\mathrm{plane}}=(3-\vp)r_{\mathrm{WZW}}^2$ and
  $(3-\vp)\cdot\sqrt{5}\vp/2=(k+2)/2$ (Lemmas~\ref{P14:lem:fmatrix_plane},
  \ref{P14:thm:higgs_algebra}, \ref{P14:thm:pentagon_g4})
\item Per-sector torus weight $= d_j^2\sin^2(\pi/(k+2))$ from KZB $b$-cycle
  monodromy (Lemma~\ref{P14:lem:kzb_monodromy}; Falceto--Gaw\c{e}dzki 1996)
\end{enumerate}
\status{proved (Theorem; all five items proved)}
\residual{$m_H = 125.10\,\mathrm{GeV}$; ATLAS $125.11\pm0.11\,\mathrm{GeV}$;
  tension $0.13\sigma$; residual $-0.01\%$.
  $\lambda=0.12907$ vs $\lambda_\mathrm{SM}=0.12909$ (residual $-0.02\%$).
  Framework $\vEW=246.32\,\mathrm{GeV}$: $m_H=125.15\,\mathrm{GeV}$ ($0.33\sigma$).}
\source{P14:thm:higgs,P14:rem:higgs_numerics,
  ATLAS PRL~\textbf{131} (2023) 251802}
\end{theorem}
\status{published}
\source{P14:thm:higgs}

\begin{remark}[Numerical verification]
\label{P14:rem:higgs_numerics}
ATLAS combined Run~1+2: $\mH=125.11\pm0.11\,\mathrm{GeV}$.
Framework: $\mH=125.10\,\mathrm{GeV}$ ($0.13\sigma$, residual $-0.01\%$).
Quartic coupling: $\lambda=0.12907$ vs $\lambda_{\mathrm{SM}}=0.12909$ (residual $-0.02\%$).
Using framework value $\vEW=246.24\,\mathrm{GeV}$: $\mH=125.11\,\mathrm{GeV}$ ($0.08\sigma$).
\status{published}
\residual{$0.13\sigma$ from ATLAS combined; $-0.01\%$}
\source{P14:rem:higgs_numerics}
\end{remark}
\status{published}
\source{P14:rem:higgs_numerics}

\begin{remark}[Algebraic structure]
\label{P14:rem:higgs_algebra}
The quartic coupling has a Fibonacci decomposition:
\begin{equation}
  \lambda = \frac{5\cdot2^{5/3}}{\vp^{10}} = \frac{5\cdot2^{5/3}}{55\vp+34},
\end{equation}
where $F_9=34$, $F_{10}=55$.
Every factor traces to the icosahedral structure: $5=k+2$ (pentagon),
$2^{5/3}=2^{(2k-1)/3}$ (WZW fusion), $\vp^{10}=(5\vp+3)^2$ (pentagon squared).
\status{published}
\source{P14:rem:higgs_algebra}
\end{remark}
\status{published}
\source{P14:rem:higgs_algebra}

\begin{remark}
\label{P14:rem:higgs_consistency}
The Higgs mass formula uses the same $r_{\mathrm{WZW}}$ that governs the toroidal
correction (Theorem~\ref{P14:thm:toroidal_formula}).
The thick-torus correction shifts $r_{\mathrm{WZW}}$ from $0.2301$ (thin) to $0.22717$ (3D),
and the Higgs mass shifts correspondingly.
The same profile ratio governs the Yukawa coupling, the electroweak VEV, and the Higgs self-coupling.
\status{published}
\source{P14:rem:higgs_consistency}
\end{remark}
\status{published}
\source{P14:rem:higgs_consistency}

\begin{remark}
\label{P14:rem:higgs_path}
The proof of Theorem~\ref{P14:thm:higgs} has five proved components:
\begin{enumerate}[noitemsep,label=(\roman*)]
\item Characters $d_j=\chi_j(e^{i\pi/(k+2)})$; $\vp=1+2\cos(2\pi/5)$
  (Lemma~\ref{P14:lem:qgroup_char}).
\item Parseval: $\sum_j\sin^2((2j+1)\pi/(k+2))=(k+2)/2$
  (Lemma~\ref{P14:lem:parseval}).
\item $D^2\sin^2(\pi/(k+2))=(k+2)/2$ (Corollary~\ref{P14:cor:D2_parseval}).
\item Plane $F$-matrix: $\lambda_{\mathrm{plane}}=(3-\vp)r_{\mathrm{WZW}}^2$
  (Lemma~\ref{P14:lem:fmatrix_plane}); $(3-\vp)\cdot\sqrt{5}\vp/2=(k+2)/2$
  (Theorem~\ref{P14:thm:higgs_algebra}).
\item KZB $b$-cycle monodromy $=$ modular S-matrix gives per-sector weight
  $d_j^2\sin^2(\pi/(k+2))$ (Lemma~\ref{P14:lem:kzb_monodromy};
  Falceto--Gaw\c{e}dzki 1996).
\end{enumerate}
Combined with Theorem~\ref{P14:thm:higgs_algebra}:
$\mH=\sqrt{k+2}\,r_{\mathrm{WZW}}\,\vEW=125.10\,\mathrm{GeV}$ derived without free parameters.
\status{published}
\source{P14:rem:higgs_path}
\end{remark}
\status{published}
\source{P14:rem:higgs_path}
